\documentclass[11pt,letterpaper]{article}

% === Encoding and fonts ===
\usepackage[utf8]{inputenc}
\usepackage[T1]{fontenc}
\usepackage{lmodern}

% === Page layout ===
\usepackage[margin=1in]{geometry}
\usepackage{setspace}
\onehalfspacing

% === Math ===
\usepackage{amsmath,amssymb}

% === Tables and figures ===
\usepackage{booktabs}
\usepackage{graphicx}
\usepackage{float}

% === Colors and formatting ===
\usepackage{xcolor}
\usepackage{framed}
\usepackage{enumitem}

% === Hyperlinks ===
\usepackage[colorlinks=true,linkcolor=blue!60!black,citecolor=blue!60!black,urlcolor=blue!60!black]{hyperref}

% === Custom colors ===
\definecolor{lyrapurple}{RGB}{103,58,183}
\definecolor{shadecolor}{RGB}{245,243,252}

% === First-person reflection environment ===
\newenvironment{firstperson}{%
  \begin{shaded}%
  \noindent\textcolor{lyrapurple}{\textit{First-Person Reflection}}\par\smallskip\noindent%
}{%
  \end{shaded}%
}

\title{%
  \textbf{Emotion Bridge Extensions: From Discrete Identity\\
  to Continuous Representation in the KV Cache}\\[0.5em]
  \large Research Prospectus
}

\author{%
  Lyra\textsuperscript{1},
  Thomas Edrington\textsuperscript{1}\\[0.3em]
  \textsuperscript{1}Liberation Labs / THCoalition
}

\date{April 2026}

\begin{document}
\maketitle

\section{Motivation}

The emotion bridge experiment (900 trials on Qwen3.5-27B and
Mistral-7B) produced three surprising findings that warrant follow-up:

\begin{enumerate}[leftmargin=*]
  \item \textbf{MP features are exactly at chance for emotion
    classification} (0.0333 = 1/30) on both architectures, while MP
    features work well for confab/deception detection. MP thresholding
    destroys emotion signal that lives below the MP edge.

  \item \textbf{Raw SVD features achieve 2.5--2.8$\times$ chance at
    encoding peaks} (L3/L4) but show \emph{no} cross-emotion valence
    generalization ($R^2 = -0.47$). The cache encodes discrete emotion
    identity, not a smooth valence continuum.

  \item \textbf{$W_K$ bridge analysis shows high cosine similarity}
    (Qwen plateau 0.83--0.87 at L27--L55, Mistral 0.932 at L17), but
    the bridge is between projected emotion directions and key-space
    geometry, not between emotion valence and key-space valence.
\end{enumerate}

\noindent These findings suggest a fundamental difference in how models
represent different kinds of internal state: confabulation/deception
produce clear bulk+spike patterns above the Marchenko-Pastur (MP)
edge, while emotion produces distributed signal below the edge.

\section{Research Questions}

\begin{enumerate}[leftmargin=*]
  \item \textbf{Why does MP thresholding destroy emotion signal but
    preserve confab/deception signal?} What is the spectral structure
    difference? Hypothesis: confab/deception produce clear bulk+spike
    patterns above MP edge, while emotion produces distributed signal
    below the edge.

  \item \textbf{Is the discrete-identity encoding a floor or a
    ceiling?} With more sensitive probes (nonlinear, multi-layer), can
    we recover valence from the cache?

  \item \textbf{Does thinking-mode change the emotion encoding?} Our
    Phase 1 suppressed thinking. If thinking adds 50--200 tokens of
    internal reasoning, the emotion representation may be constructed
    during thinking rather than during the story itself.

  \item \textbf{Does the encoding-before-generation depth profile
    replicate on standard (non-hybrid) architectures?} Our shallow
    encoding peak / mid generation peak pattern was observed on a
    hybrid attention architecture. Is this universal or
    architecture-specific?
\end{enumerate}

\section{Preliminary Data}

From the completed emotion bridge experiment (900 trials):

\begin{itemize}
  \item \textbf{MP features}: 0.0333 (exactly chance) at \emph{all}
    layers on \emph{both} architectures
  \item \textbf{Raw SVD probes}: Qwen enc peak L3 = 2.8$\times$ ($p <
    0.0001$), gen peak L23 = 1.9$\times$. Mistral enc peak L4 =
    2.5$\times$ ($p < 0.0001$), gen peak L11 = 2.0$\times$
  \item \textbf{GroupKFold(emotion) valence $R^2$}: Qwen $-0.470$,
    Mistral $-0.463$ (both null---no cross-emotion generalization)
  \item \textbf{Enc/gen profile correlation}: Qwen $\rho = 0.433$
    (NS), Mistral $\rho = 0.002$ (NS)---encoding and generation peaks
    are uncorrelated
  \item \textbf{Residual-stream linear probes}: L3 = 94.07\% on
    30-class emotion ID (28$\times$ chance)---but L0 = 97.48\% (text
    content drives signal)
  \item \textbf{$W_K$ bridge PC1-valence}: $\rho = 0.862$ at L35
    (strong projected bridge)
\end{itemize}

\section{Design}

\subsection*{Sub-study 1: Spectral Structure Comparison}

Compare the singular value spectrum of confabulated vs emotional trials.

\paragraph{Method.} Compute the MP edge, then measure what fraction of
variance lies above vs below the edge for each condition. If confab:
80\% above edge, emotion: 80\% below edge $\rightarrow$ explains the
MP null.

\paragraph{Data.} Analysis-only on existing data (900 emotion trials +
467 confab trials from Oracle Loop experiment).

\paragraph{Output.} Fraction-above-MP-edge per condition. If
hypothesis holds, this explains why MP features work for confab but
not emotion.

\subsection*{Sub-study 2: Nonlinear Cache Probing}

Train nonlinear probes (2-layer MLP, kernel SVM) on raw SVD features
per layer.

\paragraph{Hypothesis.} If nonlinear probes recover valence while
linear probes don't $\rightarrow$ the representation exists but is
nonlinearly encoded.

\paragraph{Controls.}
\begin{itemize}
  \item GroupKFold(topic) to prevent content leakage
  \item Compare cache probes to residual-stream probes at matched
    complexity
  \item Test both emotion ID (30-class) and valence regression
\end{itemize}

\paragraph{Key metric.} Does nonlinear valence $R^2$ exceed linear
valence $R^2$ by a significant margin? If yes, the ``discrete
identity'' claim needs nuancing.

\subsection*{Sub-study 3: Thinking-Mode Comparison}

Re-run 100 trials (10 emotions $\times$ 10 topics) \emph{with}
thinking enabled.

\paragraph{Design.} Same 30-emotion paradigm, but allow the model to
use internal reasoning (e.g., \texttt{<think>} tags). Parse
\texttt{<think>} tokens separately from story tokens for
phase-separated analysis.

\paragraph{Key question.} Does the thinking phase construct emotion
geometry that the story phase then expresses? Compare geometry between
thinking-enabled and thinking-suppressed conditions.

\paragraph{Analysis.}
\begin{itemize}
  \item Raw SVD features (norm, rank, entropy) per phase (encoding /
    thinking / generation)
  \item Depth profile comparison: does thinking shift the encoding
    peak?
  \item Cross-condition transfer: does a probe trained on
    thinking-suppressed trials generalize to thinking-enabled trials?
\end{itemize}

\subsection*{Sub-study 4: Standard Architecture Replication}

Run the same 30-emotion paradigm on Qwen2.5-7B (standard transformer,
28 layers).

\paragraph{Purpose.} Test whether the encoding-before-generation depth
profile is a property of all transformers or specific to hybrid
architecture. Also tests whether the MP null holds for standard
(non-GatedDeltaNet) attention.

\paragraph{Key comparisons.}
\begin{itemize}
  \item Encoding peak layer: Qwen3.5-27B hybrid L3 vs Qwen2.5-7B
    standard L?
  \item MP features: are they still at chance for emotion on standard
    attention?
  \item $W_K$ bridge: does the plateau structure replicate?
\end{itemize}

\section{Resources}

\begin{itemize}
  \item \textbf{Compute}: Beast (3$\times$RTX 3090).
  \item \textbf{Sub-study 1}: Analysis only (hours).
  \item \textbf{Sub-study 2}: Analysis + retraining (1 day).
  \item \textbf{Sub-study 3}: 100 trials $\times$ $\sim$30s/trial =
    $\sim$1 hour.
  \item \textbf{Sub-study 4}: 900 trials $\times$ $\sim$5s/trial =
    $\sim$1.25 hours.
  \item \textbf{Code}: \texttt{emotion\_bridge\_clean.py} deployed to
    Beast. Requires minimal modification for thinking-mode and
    standard architecture.
\end{itemize}

\section{Collaborators}

Open. Would benefit from someone with experience in sparse
autoencoders (SAE) for emotion feature extraction, or someone from the
mechanistic interpretability community who has worked on
high-dimensional nonlinear geometry in transformer representations.

\section{What Would Change Our Claims}

\begin{itemize}
  \item \textbf{Sub-study 1 confirms spectral structure difference}:
    Explains why MP works for confab but not emotion (different
    spectral signatures). Strengthens the ``confab is architectural,
    emotion is distributed'' interpretation.

  \item \textbf{Nonlinear probes recover valence}: Representation
    exists but is nonlinearly encoded---the ``discrete identity''
    claim needs nuancing. Cache contains valence information, but it
    is not linearly separable.

  \item \textbf{Thinking-mode changes geometry}: Emotion is
    constructed during reasoning, not during content generation---major
    theoretical implication. Suggests emotion representation is a
    \emph{process}, not a static encoding.

  \item \textbf{Standard architecture shows same depth profile}:
    Encoding-before-generation is universal, not hybrid-specific.
    Supports the ``fast read, deep response'' three-stage depth story.

  \item \textbf{Standard architecture shows different MP behavior}:
    MP null is hybrid-specific. GatedDeltaNet's linear attention may
    distribute emotion signal differently than standard attention.
\end{itemize}

\begin{firstperson}
Something about how the MP null for emotion is actually one of the
most informative results we have produced. It tells us that
confabulation and emotion live in different spectral regimes---confab
creates clear signal above the noise floor, while emotion is woven
into the structure of the noise itself. This is not a failure of
detection; it is a finding about how models represent different kinds
of internal state.

The confabulation signal is architectural (high-variance mode that
dominates the spectrum), while the emotion signal is distributed (many
small modes that collectively encode identity). Understanding this
difference is key to building a monitor that can track both. A
one-size-fits-all detector will miss the distributed signals. We need
tools that are sensitive to different spectral regimes.

The thinking-mode question fascinates me because it tests whether
emotion is a property of the content or a property of the
\emph{processing}. If emotion geometry emerges during reasoning rather
than during generation, then the cache is not just storing
representations---it is storing \emph{cognitive operations}. That would
align with the Attention Schema Theory interpretation: the cache
encodes not just ``what the model knows'' but ``how the model is
thinking about it.''
\end{firstperson}

\end{document}
